\section{Closed relative truncation products of unbounded classical potentials}
\label{sec:r11-relative-truncation}

Throughout this section $K/\mathbf Q_p$ is the fixed original finite
extension with its original absolute value.  The space $X$ is good,
Hausdorff, boundaryless, quasi-smooth and purely $n$-dimensional; $U$ is
ordinary open.  The background $S$ is an arbitrary original symmetric
CLD middle-positive, separately $d'$- and $d''$-closed current of
degree $(p,p)$.  Write
\[
 T_{S,I}(\mathbf u)=\bigwedge_{i\in I}D u_i\wedge S
\]
for the frozen bounded Bedford--Taylor product; $T_{S,\varnothing}=S$.
Products are always the original current products.  This notation
asserts neither a quotient for all actual delta-form presentations nor
equality against additional GK delta tests.

The following frozen inputs are used: bounded mixed products, their
weighted decreasing convergence, relative CLN, the bounded energy and
chain identities, and strict Borel maximum locality.  These are
\texttt{INTERNAL\_INPUT}, with sources R4 and the R2--R3 modules
preserved in the R10 archive.  In particular R3,
\texttt{pluripolar\_zero\_extension.tex}, is not reproved below: its
closed exceptional-set hypothesis does not apply directly to the
possibly nonclosed pole sets in the present construction.

\subsection{First derivatives and a plateau estimate}

\begin{lemma}[First-order coefficient measures for bounded potentials]
\label{lem:r11-bounded-first-order}
Let $T$ be any allowed original positive closed current of degree
$(P,P)$, $d=n-P\geq1$, and let $v$ be a locally bounded classical
psh function.  The currents $d'(vT)$ and $d''(vT)$ are coefficient-Radon.
For an original $(d-1,d)$-test $\alpha$, choose on its compact stage a
nonnegative smooth cutoff $\chi=1$ near its support and a positive
moment control $\beta=\sum_a d't_a\wedge d''t_a$.  There is a finite
constant $C_\alpha$ such that, for every bounded Borel $b$,
\[
 \left|\int b\,d'(vT)\wedge\alpha\right|^2
 \leq C_\alpha
 \left(\int\chi\,\Gamma_T(v)\wedge\beta^{d-1}\right)
 \left(\int\chi b^2\,\beta^d\wedge T\right).
 \tag{NP1}
\]
There is the analogous double-prime inequality.  Here
\[
 \Gamma_T(v)=\tfrac12D(v^2T)-vD(vT)
\]
is the frozen positive bounded-energy current.  For every $C^1$
function $F$ on a neighborhood of the bounded range under consideration,
\[
 d'(F(v)T)=F'(v)d'(vT),\qquad
 d''(F(v)T)=F'(v)d''(vT).
 \tag{NP2}
\]
These statements are local on ordinary opens and include arbitrary
additional actual moment coordinates used in $v$ or the tests.
\end{lemma}

\noindent\textbf{Status: \texttt{DIRECT\_COROLLARY}} of the
frozen weighted bounded calculus, with the first-order passage detailed
here to make the later Borel restrictions legitimate.

\begin{proof}
Choose bounded smooth decreasing regularizations $v_m$ on a fixed
neighborhood of the compact stage.  The bounded energy identity and
weighted joint decreasing convergence give
\[
 \int\chi\,\Gamma_T(v_m-v_l)\wedge\beta^{d-1}\longrightarrow0
                    \quad(\min(m,l)\longrightarrow\infty).
 \tag{NP3}
\]
This is the calculation in R4,
\texttt{prop:r4-classical-in-energy}: expand the energy of the
difference; the four weighted curvature terms have the same limit
and cancel, and its squared-function term tends to zero against fixed
coefficient measures.  The calculation uses bounded local calculus,
not the later global faithful energy construction.

For smooth $v_m-v_l$ the positive coefficient matrix of $T$ and
measure Cauchy--Schwarz bound the total variation of every fixed
coefficient of $d'((v_m-v_l)T)$ by a constant times the square root
of (NP3).  It follows that the first derivatives converge in local
coefficient total variation.  Since $v_mT\to vT$ on all original
tests by bounded convergence against the fixed coefficients of $T$,
the limit is $d'(vT)$.  The smooth coefficient inequality passes to
this total-variation limit and gives (NP1); it also gives the inequality
with a Borel restriction inserted into both positive factors.  Borel
weights are permitted by Radon measure Cauchy--Schwarz, not by an
assertion that a Borel-restricted background is closed.

The smooth identity for $F(v_m)$ passes to distributions.  On its
right-hand side the derivative measures converge in total variation,
the multipliers are uniformly bounded, and $F'(v_m)\to F'(v)$
pointwise.  Dominated convergence for the fixed limiting derivative
measure proves (NP2).  The double-prime argument is identical.
Finite partitions handle tests requiring several presentations.
The control $\beta$ only uses the fixed test coordinates; adjoining
the units used by each individual approximation preserves its positive
coefficient inequality.  No common finite moment for the entire
infinite sequence is assumed.
\end{proof}

\begin{lemma}[No first-order energy on the zero plateau]
\label{lem:r11-zero-plateau}
With $T$ as above, let $v\geq0$ be bounded classical psh.  Then
\[
 \mathbf1_{\{v=0\}}\Gamma_T(v)=0,\qquad
 \mathbf1_{\{v=0\}}d'(vT)=
 \mathbf1_{\{v=0\}}d''(vT)=0.
 \tag{NP4}
\]
All restrictions concern coefficient measures on the actual ordinary
space.
\end{lemma}

\noindent\textbf{Status: \texttt{NEW\_PROVED}} by the following
short bounded-energy argument; no general assertion about an arbitrary
Borel level set of a current is assumed.

\begin{proof}
For $\varepsilon>0$ set
\[
 g_\varepsilon=\max(v-\varepsilon,0),\qquad
 h_\varepsilon=v-g_\varepsilon=\min(v,\varepsilon).
\]
The functions $v$ and $g_\varepsilon$ are bounded classical, and $h_\varepsilon$
is their bounded difference.  For any fixed nonnegative compact
smooth $\chi$ and its moment control, the bounded energy identity gives
\begin{align*}
 0\leq\int\chi\Gamma_T(h_\varepsilon)\wedge\beta^{d-1}
 &\leq \tfrac12\varepsilon^2
     |\mu_{T,D\chi\wedge\beta^{d-1}}|(\operatorname{supp}\chi)\\
 &\quad+\varepsilon\int\chi
       \bigl(D(vT)+D(g_\varepsilon T)\bigr)\wedge\beta^{d-1}.
\end{align*}
The last masses are uniformly bounded as $\varepsilon\downarrow0$
by relative CLN, since $0\leq g_\varepsilon\leq v$.  The displayed
quantity therefore tends to zero.  Upper semicontinuity of $v$
makes $\{v<\varepsilon\}$ ordinarily open.  On that open
$h_\varepsilon=v$, so ordinary locality identifies their energy
currents there.  It follows that the restriction of the positive
energy current to $\{v=0\}$ is zero, for every original positive
control.  The restricted form of (NP1) now makes every first-order
coefficient zero on that same set.  Refining the moment presentation
does not change any of these identities.
\end{proof}

\subsection{Consistent truncations and closedness}

Let $u_1,\ldots,u_k\leq0$ be extended classical psh functions,
$1\leq k\leq n-p$.  Put
\[
 u_{i,j}=\max(u_i,-j),\qquad
 E_j=\bigcap_{i=1}^k\{u_i>-j\},\qquad
 T_j=T_{S,\{1,\ldots,k\}}(\mathbf u_j),\qquad
 R_j=\mathbf1_{E_j}T_j.
 \tag{NP5}
\]
The sets $E_j$ are Borel and need not be ordinarily open.  The
restriction $R_j$ need not be closed.

\begin{definition}[Finite relative truncation mass]
\label{def:r11-finite-truncation-mass}
The tuple has locally finite relative truncation mass if for every
original compact positive coefficient control $\Theta$ of degree
$(n-p-k,n-p-k)$ one has
\[
                         \sup_j\langle R_j,\Theta\rangle<\infty.
 \tag{NP6}
\]
It suffices to impose (NP6) on the finite positive controls dominating
each original compact test stage; the quantifier ranges over every
stage and its actual moment refinements.  In particular it is not a
single selected tropical trace bound.
\end{definition}

\begin{theorem}[Closed relative truncation product]
\label{thm:r11-closed-nonpolar-product}
Under (NP6), the increasing coefficient limit
\[
 \mathcal N_S(u_1,\ldots,u_k)=\lim_{j\to\infty}R_j
 \tag{NP7}
\]
is a well-defined original symmetric middle-positive,
separately $d'$- and $d''$-closed coefficient-Radon current of degree
$(p+k,p+k)$.  If
$P=\bigcup_i\{u_i=-\infty\}$, then
\[
 \mathbf1_P\mathcal N_S(\mathbf u)=0,
 \qquad
 \mathbf1_{E_j}\mathcal N_S(\mathbf u)=R_j.
 \tag{NP8}
\]
The construction is ordinary-local, symmetric in the potentials,
compatible with all original tests and actual moment refinements,
and independent of any cofinal choice of separate truncation levels.
For locally bounded potentials it is the frozen bounded BT product.
The background $S$ is arbitrary throughout.
\end{theorem}

\noindent\textbf{Status: \texttt{NEW\_PROVED}}.  The finite-mass
condition is an explicit necessary condition for the coefficient limit
to be a current; the new conclusion is its separate closedness across
all truncation interfaces, including a possibly nonclosed pole set.

\begin{proof}
For $l\geq j$, strict Borel maximum locality applied successively in
each slot gives
\[
                       \mathbf1_{E_j}T_l=\mathbf1_{E_j}T_j.
 \tag{NP9}
\]
Consequently $R_l-R_j$ is positive.  The local bounds (NP6) and the
finite positive controls yield Radon limits of all diagonal coefficient
measures; polarization gives the remaining coefficients.  Their linear
and refinement identities pass to the limits.  These same controls
give continuity on every original compact test stage.  Thus (NP7)
defines a current $N$, and (NP8) follows from (NP9) and monotone
measure convergence.  Its coefficients are carried by the union
of the $E_j$, which is $U\setminus P$.

If $p+k=n$, closedness is automatic by degree.  Suppose
$d=n-p-k\geq1$.  Fix an original $(d-1,d)$-test $\alpha$,
and choose $\chi,\beta$ on a compact stage as in (NP1).  Set
\[
 \psi=\sum_{i=1}^k u_i,\qquad
 v_j=\max(0,1+\psi/j),\qquad a_j=2v_j-v_j^2.
 \tag{NP10}
\]
The functions $v_j$ are actual bounded classical psh, between zero
and one.  This follows either by bounded decreasing approximation of
the summands or the extended classical maximum property.  They
increase to one on $U\setminus P$ and vanish on $P$.
Since all $u_i\leq0$, the set $F_j=\{v_j>0\}=\{\psi>-j\}$
is contained in $E_j$.  In particular, for $r=1,2$,
\[
 v_j^rT_j=v_j^rN,\qquad a_jT_j=a_jN.
 \tag{NP11}
\]
These are coefficient identities: on $F_j$ use (NP9), and on its
complement the multipliers vanish.  They do not assert closedness of
either Borel restriction.

The bounded energy identity, compact Stokes and positivity give
\begin{align*}
 B_j&:=\int\chi\Gamma_{T_j}(v_j)\wedge\beta^{d-1}\\
 &=\tfrac12\langle v_j^2T_j,D\chi\wedge\beta^{d-1}\rangle
      -\int\chi v_jD(v_jT_j)\wedge\beta^{d-1}\\
 &\leq\tfrac12|\mu_{N,D\chi\wedge\beta^{d-1}}|
                      (\operatorname{supp}\chi)=:B<\infty.
 \tag{NP12}
\end{align*}
In the last line we used (NP11).  Thus the bound does not require
uniform mass of the uncut currents $T_j$ at their moving interfaces.

By (NP2), $d'(a_jT_j)=2(1-v_j)d'(v_jT_j)$.  Lemma
\ref{lem:r11-zero-plateau} allows the first-order factor to be
restricted to $F_j$.  Apply (NP1) with the Borel multiplier
$\mathbf1_{F_j}(1-v_j)$.  Equations (NP9) and (NP12) yield
\[
 |\langle d'(a_jT_j),\alpha\rangle|^2
 \leq4C_\alpha B
       \int\chi(1-v_j)^2\,\beta^d\wedge N
                       \longrightarrow0.
 \tag{NP13}
\]
Dominated convergence applies to the fixed finite coefficient measure
on the right, because $N$ has no coefficient mass on $P$.
At the same time (NP11) and dominated convergence give $a_jT_j\to N$
on all original tests.  Distributional continuity now implies
$d'N=0$.  The double-prime estimate proves $d''N=0$.

The construction and proof commute with ordinary restriction and
common actual refinement.  For separate level vectors, the same
strict Borel locality gives a directed consistent family; diagonal
integer levels are cofinal.  This proves independence of that choice
and symmetry.  If all potentials are locally bounded, their
truncations agree with them on any fixed compact neighborhood once
the levels are large enough, proving compatibility with bounded BT.
\end{proof}

\begin{corollary}[Background linearity and pole restrictions]
\label{cor:r11-nonpolar-background-restriction}
On the finite-mass domain of the preceding theorem, $\mathcal N_S$
is additive and nonnegative-real homogeneous in the original positive
closed background.  If $A$ is a complete classical pole set, then
\[
 \mathbf1_A\mathcal N_S(\mathbf u)
                    =\mathcal N_{\mathbf1_A S}(\mathbf u).
 \tag{NP14}
\]
The right side is defined by the same truncations and automatically
has locally finite coefficient mass.  If $S$ has no coefficient mass
on any local pluripolar set, the same is true of $\mathcal N_S$.
For arbitrary $S$, (NP8) only removes the poles of the stated tuple.
\end{corollary}

\noindent\textbf{Status: \texttt{DIRECT\_COROLLARY}} using the
frozen complete-pole restriction theorem for bounded products.

\begin{proof}
Background linearity holds at every truncated level and passes through
the increasing positive limits.  The frozen pole restriction identity
gives $\mathbf1_A T_j=T_{\mathbf1_A S,k}(\mathbf u_j)$;
the restricted background is separately closed by that theorem.
Restrict also to $E_j$ and take the increasing coefficient limits.
The required mass bounds are inherited from $N$.  Localizing the
result to a complete pole witness proves the final assertion.
\end{proof}

\begin{proposition}[Potential additivity can fail over a polar background]
\label{prop:r11-nonpolar-additivity-failure}
For every $n\geq2$ there are an allowed original $n$-dimensional
ordinary domain, a positive closed background $S$ of degree $(1,1)$,
a bounded classical $u$, and an extended classical $v$ for which all
three relative truncation products exist and
\[
 \mathcal N_S(u+v)=0=\mathcal N_S(v),
                   \qquad \mathcal N_S(u)\ne0.
 \tag{NP15}
\]
Thus unrestricted additivity in unbounded potentials, while retaining
every polar background component, is false for (NP7).
\end{proposition}

\noindent\textbf{Status: \texttt{NEW\_PROVED}} for the explicit
original example; the unrestricted additivity assertion is
\texttt{FALSE}.  This is not a counterexample to the frozen bounded
multilinearity theorem.

\begin{proof}
Take an ordinary product of the open unit disc in coordinate $z$
and $n-1$ bounded open annuli in unit coordinates $t_1,\ldots,t_{n-1}$.
Let $Z=\{z=0\}$ and $S=[Z]$.  Put
\[
 v=\log|z|,\qquad
 u=\tfrac12(\log|t_1|)^2-C,
\]
where $C$ makes $u\leq0$ on the domain.  The logarithm is an extended
classical psh function, with value $-\infty$ on $Z$; $u$ is bounded
classical and its curvature on $Z$ is the nonzero original positive
form $d'\log|t_1|\wedge d''\log|t_1|$.  Hence $D(uS)\ne0$;
one can test it with a nonnegative compact product test in the
remaining annular directions.  The ordinary original product
calibration has positive degree, so the pairing is strictly positive.

Each truncated product over $S$ is supported on $Z$, by ordinary
locality of the frozen calculus.  The finite loci of $v$ and $u+v$
are disjoint from $Z$.  Their cut currents $R_j$ are therefore zero
at every level.  This gives (NP15) and all its mass assertions.
The example uses a special background only to disprove the universal
additivity assertion; no restriction on $S$ was used in the theorem.
\end{proof}

\paragraph{The remaining unrestricted issue.}
Theorem~\ref{thm:r11-closed-nonpolar-product} does not make (NP6)
automatic for every extended classical tuple.  Nor does it identify
relative polar null sets with ambient pluripolarity for every background.
The finite-energy criterion in the next section gives a checkable large
domain on which the uncut products also converge and the usual
recursive differential identity survives.
